\section{目标函数与噪声模型的对应：从生成机制到优化目标}
\label{sec:13-Machine-Learning/01-Learning-Problems-and-Evaluation/03}

为什么回归用 MSE、分类用交叉熵？这些选择来自数据的噪声结构、训练目标和正则化约束。似然项对应噪声模型，正则项对应先验约束，两者合起来才是完整的优化目标。

\subsection{先感受困难：一个异常值毁掉整个回归}

考虑简单线性回归 \(y_i=\beta x_i+\varepsilon_i\)，用 MSE 估计 \(\beta\)。若数据干净，MSE 给出最优估计；但若混入一个异常值 \((x_{n+1},y_{n+1})=(100,1000)\)，MSE 会把它当作``信息量最大的点''——残差平方 \((y-\beta x)^2\) 随距离平方增长，异常值主导目标函数，把 \(\widehat\beta\) 拉向自己。

MAE（绝对损失）对这个异常值不敏感：残差线性增长，异常值的影响与它的距离成正比，不会被平方放大。但 MAE 在干净数据上效率低于 MSE——对 Gaussian 噪声，MSE 是最优的（Cramér--Rao 下界），MAE 只是鲁棒的折中。

没有普适最优的损失，只有与噪声模型匹配的损失。MSE 对 Gaussian 噪声最优，MAE 对 Laplace 噪声最优，Huber 对 Gaussian+异常值混合最优。选择损失函数，就是选择对噪声的假设。

\subsection{似然项：噪声模型决定目标函数的核心}

噪声模型决定目标函数的机制是最大似然。设 \(Y\given X=x\) 的条件分布为 \(p(y\given x;\theta)\)，对独立样本 \((x_i,y_i)\)，似然为

\[
L(\theta)=\prod_{i=1}^n p(y_i\given x_i;\theta).
\]

最大化似然等价于最小化负对数似然

\[
-\log L(\theta)=-\sum_{i=1}^n\log p(y_i\given x_i;\theta).
\]

不同的噪声模型给出不同的负对数似然。若噪声是 Gaussian 的，即 \(Y=\mu(x)+\varepsilon\)，\(\varepsilon\sim\mathcal N(0,\sigma^2)\)，负对数似然是

\[
\frac{1}{2\sigma^2}\sum_{i=1}^n(y_i-\mu(x_i))^2+\text{常数},
\]

这正是 MSE（差一个常数因子）。把噪声换成 Bernoulli，\(Y\sim\operatorname{Bernoulli}(p(x))\)，负对数似然变成

\[
-\sum_{i=1}^n\bigl[y_i\log p(x_i)+(1-y_i)\log(1-p(x_i))\bigr],
\]

这正是交叉熵。若噪声是 Laplace 的，\(\varepsilon\sim\operatorname{Laplace}(0,b)\)，负对数似然是 \(\frac1b\sum\abs{y_i-\mu(x_i)}+\text{常数}\)，即 MAE。若噪声服从 Poisson，\(Y\sim\operatorname{Poisson}(\lambda(x))\)，负对数似然是 \(\sum[\lambda(x_i)-y_i\log\lambda(x_i)]\)，即 Poisson 损失。

选择 MSE 就是假设 Gaussian 噪声，选择交叉熵就是假设 Bernoulli 噪声。这个对应不是查表，而是最大似然的直接结果：给定生成模型，似然项自然导出损失形式。

\subsection{训练目标：似然项只是起点}

似然项由噪声模型决定，但训练目标可以更进一步：预测均值、中位数、分位数、概率还是排序，由任务需求决定。

MSE 对应条件均值 \(\E[Y\given X]\)，在 Gaussian 噪声下最优。MAE 对应条件中位数，在 Laplace 噪声下最优，对异常值鲁棒。分位数损失（pinball loss）对应条件分位数，是不对称代价（如缺货 vs 报废）下的最优选择。

当任务要的是概率而非点预测时，交叉熵对应条件概率 \(P(Y\given X)\)，在 Bernoulli 噪声下最优，proper scoring rule 保证校准。当绝对数值本身不重要时，pairwise loss（如 RankNet）对应相对顺序，不关心绝对值，只关心排序。

同样的 Gaussian 噪声下，可以选择预测均值（MSE）、中位数（MAE）或分位数（pinball loss）。似然项给出与噪声匹配的目标，训练目标给出与任务匹配的目标。

\subsection{正则化项：先验约束优化结果}

似然项拟合数据，正则化项约束模型。完整的目标函数是

\[
J(\theta)=-\log L(\theta)+\lambda\Omega(\theta),
\]

其中 \(\Omega(\theta)\) 是正则化项，\(\lambda\) 控制强度。正则化把领域知识编码进优化。

L2（Ridge）取 \(\Omega(\theta)=\norm{\theta}_2^2\)，假设权重小（平滑），改善病态方向，对应 Gaussian 先验。L1（Lasso）取 \(\Omega(\theta)=\norm{\theta}_1\)，假设权重稀疏（特征选择），产生零权重，对应 Laplace 先验。Elastic Net 取 \(\Omega(\theta)=\alpha\norm{\theta}_1+(1-\alpha)\norm{\theta}_2^2\)，兼顾稀疏与平滑。

并非所有正则化都写成惩罚项。早停在训练到验证误差最小时停止，是隐式正则化。Dropout 在训练时随机失活神经元，假设特征不应共适应。

MAP（最大后验）视角统一了似然与正则化：

\[
\theta_{\text{MAP}}=\argmax_\theta\log p(y\given x;\theta)+\log p(\theta),
\]

其中 \(\log p(\theta)\) 是先验。Gaussian 先验对应 L2，Laplace 先验对应 L1。

\subsection{三个维度的统一：噪声、任务、约束}

目标函数的选择是三个维度的平衡：噪声模型决定似然项的数学形式（Gaussian→MSE，Bernoulli→交叉熵，Laplace→MAE）；任务需求决定预测目标（均值、中位数、概率还是排序）；正则化约束模型容量（L2 平滑、L1 稀疏、早停、Dropout）。

三个维度可以独立变化：Gaussian 噪声下可以选择预测均值（MSE）或中位数（MAE），可以加 L2 或 L1 正则化。但维度之间也有交互：强正则化可能偏离似然项的最优解（偏差-方差权衡），任务目标可能要求特定的似然形式（分位数预测需要分位数损失，不能用 MSE 替代）。

\subsection{目标函数的形状决定优化几何}

目标函数的形状——凸性、光滑性、增长速率——决定优化的难度与性质。

MSE 是二次的：梯度线性，Hessian 常数，构成凸优化问题，线性回归有闭式解（正规方程）；但二次增长对异常值敏感——残差平方随距离平方增长。交叉熵是凸的：梯度含 sigmoid/softmax，需要迭代优化（梯度下降、Newton）；凸性保证全局最优，但条件数可能很差（类别不平衡时）。

MAE 是分段线性的，只有次梯度（在零点不可导），对异常值鲁棒；但线性增长意味着对 Gaussian 噪声效率低，渐近方差比 MSE 大 \(\pi/2\) 倍。Hinge 同样是分段线性的，只有次梯度；它给出支持向量（稀疏解），不可导点对应决策边界，次梯度提供下降方向。

形状决定优化几何：二次损失给出光滑凸问题，分段线性损失给出非光滑凸问题，0-1 损失给出非凸组合问题。

\subsection{目标函数的统计性质决定泛化}

目标函数的统计性质——偏差-方差、校准、proper scoring rule——决定泛化能力。

MSE 对 Gaussian 噪声有效（达到 Cramér--Rao 下界），但对异常值高方差；MAE 对异常值低方差，但对 Gaussian 噪声低效。交叉熵是 strictly proper scoring rule——最优预测是真实条件概率 \(p(y\given x)\)，不是``最能分开两类的分数''。

0-1 损失是最优的分类目标（Bayes 风险），但不可导。Hinge、logistic、交叉熵都是凸替代——在一定条件下（如 \(p(y\given x)\) 单调），替代损失的最优解与 0-1 损失一致，但替代损失本身可能过度惩罚``容易分类的点''（logistic）或``边界附近的点''（hinge）。

\subsection{这些对应关系都有前提}

对应关系成立的前提是噪声模型正确：目标函数与噪声模型匹配时最优，错配时失效。MSE 对 Gaussian 噪声最优，但对 Laplace 噪声低效；交叉熵对 Bernoulli 噪声最优，但对重尾噪声（如 t 分布）可能过度自信。若真实噪声是混合分布，任何单一噪声假设的目标函数都只是近似；鲁棒方法（Huber、trimmed loss）是折中。

即使噪声模型正确，目标函数里仍有需要外部选择的量。正则化强度 \(\lambda\) 是超参数，需要交叉验证选择：过强的正则化导致欠拟合，过弱的正则化导致过拟合。替代损失与真实损失之间也有差距：0-1 损失不可导，凸替代（hinge、logistic）在一定条件下一致，但替代损失本身可能过度惩罚某些点。

目标函数还与优化算法交互：MSE 加线性模型是闭式解，MSE 加深度网络是非凸优化；交叉熵加线性模型是凸优化，交叉熵加深度网络是非凸优化。

\subsection{落点}

目标函数的选择是三个维度的统一：似然项匹配噪声模型（Gaussian→MSE，Bernoulli→交叉熵），训练目标超越噪声（预测均值、中位数、概率、排序），正则化项约束模型（L2 平滑、L1 稀疏、早停、Dropout）。噪声模型给出似然，先验给出正则化，任务需求选择预测目标。目标函数的形状决定优化几何，统计性质决定泛化——选择目标函数就是选择任务定义。
